
\documentclass{article}
\usepackage{colortbl}
\usepackage{makecell}
\usepackage{multirow}
\usepackage{supertabular}

\begin{document}

\newcounter{utterance}

\centering \large Interaction Transcript for game `bbh', experiment `formal\_fallacies', episode 223 with Qwen3.6{-}27B.
\vspace{24pt}

{ \footnotesize  \setcounter{utterance}{1}
\setlength{\tabcolsep}{0pt}
\begin{supertabular}{c@{$\;$}|p{.15\linewidth}@{}p{.15\linewidth}p{.15\linewidth}p{.15\linewidth}p{.15\linewidth}p{.15\linewidth}}
    \# & \multicolumn{2}{c}{Player} && \multicolumn{2}{c}{Game Master} \\
    \hline

    \theutterance \stepcounter{utterance}  
    & & & \multicolumn{4}{p{0.6\linewidth}}{
        \cellcolor[rgb]{0.9,0.9,0.9}{
            \makecell[{{p{\linewidth}}}]{
                \texttt{\tiny{[P1$\langle$GM]}}
                \texttt{Distinguish deductively valid arguments from formal fallacies.} \\
\\ 
\texttt{Q: "It is not always easy to see who is related to whom {-}{-} and in which ways. The following argument pertains to this question: To begin with, Lesley is a close friend of Fernando. Moreover, being a close friend of Fernando or a schoolmate of Lowell is sufficient for being a great{-}grandfather of Leroy. It follows that Lesley is a great{-}grandfather of Leroy."} \\
\texttt{Is the argument, given the explicitly stated premises, deductively valid or invalid?} \\
\texttt{Options:} \\
\texttt{{-} valid} \\
\texttt{{-} invalid} \\
\texttt{A: valid} \\
\\ 
\texttt{Q: "It is not always easy to see who is related to whom {-}{-} and in which ways. The following argument pertains to this question: Whoever is not a great{-}grandfather of Clyde is a stepbrother of Brian. Being an ancestor of Dana is sufficient for not being a great{-}grandfather of Clyde. We may conclude: Everyone who is an ancestor of Dana is a stepbrother of Brian, too."} \\
\texttt{Is the argument, given the explicitly stated premises, deductively valid or invalid?} \\
\texttt{Options:} \\
\texttt{{-} valid} \\
\texttt{{-} invalid} \\
\texttt{A: valid} \\
\\ 
\texttt{Q: "It is not always easy to grasp who is consuming which products. The following argument pertains to this question: Every infrequent user of Paul Mitchell shampoo is either a rare consumer of Nioxin shampoo or a loyal buyer of Caress soap, or both. No regular consumer of Lush soap is a rare consumer of Nioxin shampoo and, in the same time, a loyal buyer of Caress soap. It follows that whoever is an infrequent user of Paul Mitchell shampoo is not a regular consumer of Lush soap."} \\
\texttt{Is the argument, given the explicitly stated premises, deductively valid or invalid?} \\
\texttt{Options:} \\
\texttt{{-} valid} \\
\texttt{{-} invalid} \\
\texttt{A: invalid} \\
\\ 
\texttt{Q: "Is Fred a fan of Liverpool? Are supporters of Real Madrid devotees of PSG? In European football, it is sometimes difficult to keep track of the mutual admiration and dislike. The following argument seeks to clarify some such relations: First of all, no expert of Malmö FF who is an ex{-}supporter of Brøndby IF is an ex{-}fan of F91 Dudelange. Next, whoever is not a follower of Stade Rennais FC is an ex{-}fan of F91 Dudelange. So, necessarily, whoever is none of this: an expert of Malmö FF or ex{-}supporter of Brøndby IF, is a follower of Stade Rennais FC."} \\
\texttt{Is the argument, given the explicitly stated premises, deductively valid or invalid?} \\
\texttt{Options:} \\
\texttt{{-} valid } \\
\texttt{{-} invalid} \\
\texttt{A:} \\
            }
        }
    }
     \\ \\

    \theutterance \stepcounter{utterance}  
    & \multicolumn{4}{p{0.6\linewidth}}{
        \cellcolor[rgb]{0.561,0.737,0.561}{
            \makecell[{{p{\linewidth}}}]{
                \texttt{\tiny{[P1$\rangle$GM]}}
                \texttt{To determine the validity of the argument, let's translate the statements into logical notation.} \\
\\ 
\texttt{**1. Define the predicates:**} \\
\texttt{Let \$M(x)\$: \$x\$ is an expert of Malmö FF.} \\
\texttt{Let \$B(x)\$: \$x\$ is an ex{-}supporter of Brøndby IF.} \\
\texttt{Let \$D(x)\$: \$x\$ is an ex{-}fan of F91 Dudelange.} \\
\texttt{Let \$S(x)\$: \$x\$ is a follower of Stade Rennais FC.} \\
\\ 
\texttt{**2. Translate the premises:**} \\
\texttt{*   **Premise 1:** "No expert of Malmö FF who is an ex{-}supporter of Brøndby IF is an ex{-}fan of F91 Dudelange."} \\
\texttt{    *   This means if someone is both \$M\$ and \$B\$, they are not \$D\$.} \\
\texttt{    *   Logic: \$\textbackslash{}forall x ((M(x) \textbackslash{}land B(x)) \textbackslash{}rightarrow \textbackslash{}neg D(x))\$} \\
\texttt{    *   Contrapositive: \$\textbackslash{}forall x (D(x) \textbackslash{}rightarrow \textbackslash{}neg (M(x) \textbackslash{}land B(x)))\$ which is equivalent to \$\textbackslash{}forall x (D(x) \textbackslash{}rightarrow (\textbackslash{}neg M(x) \textbackslash{}lor \textbackslash{}neg B(x)))\$.} \\
\\ 
\texttt{*   **Premise 2:** "Whoever is not a follower of Stade Rennais FC is an ex{-}fan of F91 Dudelange."} \\
\texttt{    *   Logic: \$\textbackslash{}forall} \\
            }
        }
    }
    & & \\ \\

    \theutterance \stepcounter{utterance}  
    & & & \multicolumn{2}{p{0.3\linewidth}}{
        \cellcolor[rgb]{0.9,0.9,0.9}{
            \makecell[{{p{\linewidth}}}]{
                \texttt{\tiny{[GM$|$GM]}}
                \texttt{Target: invalid} \\
            }
        }
    }
    & & \\ \\

    \theutterance \stepcounter{utterance}  
    & & & \multicolumn{2}{p{0.3\linewidth}}{
        \cellcolor[rgb]{0.9,0.9,0.9}{
            \makecell[{{p{\linewidth}}}]{
                \texttt{\tiny{[GM$|$GM]}}
                \texttt{game\_result = LOSE} \\
            }
        }
    }
    & & \\ \\

\end{supertabular}
}

\end{document}
